%%%%%%%%%%%%%%%%%%%%%%%%%%%%%%%%%%%%%%%%%%%%%%%%%%%%%%%%%%%%%%%%%%%%%%
% Lớp 8 | Chương III. Tứ giác | Hình thang cân
% Dạng 8H3?11-1: Tính chất hình thang, hình thang cân
%%%%%%%%%%%%%%%%%%%%%%%%%%%%%%%%%%%%%%%%%%%%%%%%%%%%%%%%%%%%%%%%%%%%%%


% nguồn: De_AI_Toan8_cuoi_hoc_ki_i_1 (nhập 21/11/2026)
\begin{bt}%[8H3H11-1]%[ma:88e09bce]
	Cho hình thang cân $ABCD$ ($AB \parallel CD$, $AB < CD$). Gọi $O$ là giao điểm của hai đường chéo $AC$ và $BD$.
a) Chứng minh tam giác $ACD$ bằng tam giác $BDC$.
b) Chứng minh $OA = OB$ và $OC = OD$.
	\loigiai{
		a) Vì $ABCD$ là hình thang cân ($AB \parallel CD$) nên $AD = BC$ và $AC = BD$. \\
Xét $\triangle ACD$ và $\triangle BDC$ có: $AD = BC$, $AC = BD$, $CD$ chung. Suy ra $\triangle ACD = \triangle BDC$ (c.c.c). \\
b) Từ $\triangle ACD = \triangle BDC$ suy ra $\widehat{ACD} = \widehat{BDC}$, hay $\widehat{OCD} = \widehat{ODC}$. \\
Suy ra tam giác $OCD$ cân tại $O$, do đó $OC = OD$. \\
Mặt khác $AC = BD$, nên $AC - OC = BD - OD \Rightarrow OA = OB$.
	}
\end{bt}

% nguồn: De_AI_Toan8_cuoi_hoc_ki_i_1 (nhập 21/11/2026)
\begin{bt}%[8H3H11-1]%[ma:9ec4af2f]
	Cho hình thang cân $ABCD$ ($AB \parallel CD$) có $AB = AD = BC = 6\text{ cm}$ và $\widehat{D} = 60^\circ$.
a) Tính độ dài đáy lớn $CD$.
b) Chứng minh rằng $AC \perp AD$.
	\loigiai{
		a) Kẻ hai đường cao $AH \perp CD$, $BK \perp CD$ ($H, K \in CD$). Ta có $HK = AB = 6\text{ cm}$ và $DH = CK$. \\
Tam giác $AHD$ vuông tại $H$ có $\widehat{D} = 60^\circ \Rightarrow \widehat{DAH} = 30^\circ$, nên $DH = \dfrac{1}{2}AD = 3\text{ cm}$. \\
Suy ra $CK = 3\text{ cm}$ và $CD = DH + HK + CK = 3 + 6 + 3 = 12\text{ cm}$. \\
b) Vì $AB = BC$ nên tam giác $ABC$ cân tại $B$, suy ra $\widehat{BAC} = \widehat{BCA}$. \\
Vì $AB \parallel CD$ nên $\widehat{BAC} = \widehat{ACD}$ (so le trong), suy ra $\widehat{BCA} = \widehat{ACD}$, tức $CA$ là phân giác góc $BCD$. \\
Mà $\widehat{BCD} = \widehat{D} = 60^\circ$ nên $\widehat{ACD} = \dfrac{60^\circ}{2} = 30^\circ$. \\
Xét tam giác $ACD$ có $\widehat{DAC} = 180^\circ - (\widehat{D} + \widehat{ACD}) = 180^\circ - (60^\circ + 30^\circ) = 90^\circ$. Vậy $AC \perp AD$.
	}
\end{bt}
